\documentclass[a4paper,landscape,headrule,footrule]{foils}
%\usepackage{nttfoilhead}
%\newcommand{\myslide}[1]{\foilhead[-25mm]{\raisebox{12mm}[0mm]{\emp{#1}}}}
%\newcommand{\myslider}[1]{\rotatefoilhead[-25mm]{\raisebox{12mm}[0mm]{\emp{#1}}}}
%\newcommand{\myslider}[1]{\rotatefoilhead{\raisebox{-8mm}{\emp{#1}}}}
%%
%% FIXME: not all of the features are smallcaps :-(
%%
%%
%%% macros for Theories of Grammar
%%%
\usepackage{polyglossia}
\setdefaultlanguage{english}
%\setmainfont{TeX Gyre Pagella}


\newcommand{\logo}{~}
\newcommand{\header}[3]{%
\title{\vspace*{-2ex} \large JMORF/MORS --- Morpho-Syntax
\\[2ex] \Large  \emp{#2} \\ \emp{#3}}
\author{\blu{Francis Bond}   \\ 
\normalsize  \textbf{Palacký University}\\
\normalsize  \url{https://fcbond.github.io/}\\
\normalsize  \texttt{bond@ieee.org}}
\MyLogo{JMORF (2026)}
\renewcommand{\logo}{#2}
\hypersetup{
   pdfinfo={
     Author={Francis Bond},
     Title={#1: #2},
     Subject={JMORF/MORS: Morpho-Syntax},
     Keywords={Syntax, Morphology, HPSG, Unification, Constructions},
     License={CC BY 4.0}
   }
 }

\date{#1
  \\ Location: KB-1.11}
}
\usepackage{hyperref}
\hypersetup{
     colorlinks,
     linkcolor={blue!50!black},
     citecolor={red!50!black},
     urlcolor={blue!80!black}
   }



\usepackage{xcolor}
\usepackage{graphicx}
\newcommand{\blu}[1]{\textcolor{blue}{#1}}
\newcommand{\grn}[1]{\textcolor{green}{#1}}
\newcommand{\hide}[1]{\textcolor{white}{#1}}
\newcommand{\emp}[1]{\textcolor{red}{#1}}
\newcommand{\txx}[1]{\textbf{\textcolor{blue}{#1}}}
\newcommand{\lex}[1]{\textbf{\mtcitestyle{#1}}}


% \renewcommand{\labelitemi}{\textcolor{violet}{‣}}
% \renewcommand{\labelitemii}{\textcolor{purple}{‣}}

\newcommand{\subhead}[1]{\noindent\textbf{#1}\\[5mm]}

\newcommand{\Bad}{\emp{\raisebox{0.15ex}{\ensuremath{\mathbf{\otimes}}}}}
\newcommand{\bad}[1]{*\eng{#1}}

\newcommand{\com}[1]{\hfill (\emp{#1})}%

\usepackage{relsize,xspace}
\newcommand{\into}{\ensuremath{\rightarrow}\xspace}
\newcommand{\tot}{\ensuremath{\leftrightarrow}\xspace}
\usepackage{url}
\newcommand{\lurl}[1]{\MyLogo{\url{#1}}}

\usepackage{langsci-gb4e}
\let\eachwordone=\slshape
%%% langsci-gb4e defines \thexnumi as \@roman\@xsi{xnumi}, which only works
%%% inside an exe environment; some decks use xlisti on its own.  Restore
%%% the plain gb4e definition so that keeps working.
\makeatletter
\def\thexnumi{\@xsi{xnumi}}
\makeatother
\newcommand{\lx}[1]{\textbf{\mtciteform{#1}}}
\newcommand{\ix}{\ex\slshape}



\newcommand{\ent}{\ensuremath{\Rightarrow}\xspace}
\newcommand{\ngv}{\ensuremath{\not\Rightarrow}\xspace}
%\usepackage{times}
%\usepackage{nttfoilhead}
\newcommand{\myslide}[1]{\foilhead[-25mm]{\raisebox{12mm}[0mm]{\emp{#1}}}\MyLogo{\logo}}
\newcommand{\myslider}[1]{\rotatefoilhead[-25mm]{\raisebox{12mm}[0mm]{\emp{#1}}}}
%\newcommand{\myslider}[1]{\rotatefoilhead{\raisebox{-8mm}{\emp{#1}}}}

\newcommand{\section}[1]{\myslide{}{\begin{center}\Huge \emp{#1}\end{center}}}



\usepackage[lyons,j,e,k]{mtg2e}
%\renewcommand{\mtcitestyle}[1]{\textcolor{-red!75!green!50}{\textsl{#1}}}
\renewcommand{\mtcitestyle}[1]{\textcolor{teal}{\textsl{#1}}}
\newcommand{\cs}{\mtciteform}
\newcommand{\iz}[1]{\texttt{\textup{#1}}}
\newcommand{\gm}{\textsc}
\usepackage[normalem]{ulem}
\newcommand{\ul}{\uline}
\newcommand{\ull}{\uuline}
\newcommand{\wl}{\uwave}
\newcommand{\vs}{\ensuremath{\Leftrightarrow}~}
%%%
%%% Bibliography
%%%
\usepackage{natbib}
%\usepackage{url}
\usepackage{bibentry}



\usepackage[utf8]{inputenc}
%%% trees
\usepackage{multicol}
%%% pst-node: one example in lec-12-control draws a connector between
%%% two points with \rnode/\ncbar.  This used to come in with rtrees.
\usepackage{pst-node}
\usepackage{forest}
\useforestlibrary{linguistics}
\forestapplylibrarydefaults{linguistics}

%%% \makebox's optional arguments are read as AVM delimiters by
%%% langsci-avm, so hide them behind a macro with no visible brackets.
\newcommand{\lbox}[1]{\makebox[0mm][l]{#1}}

%%% Tree styles.
%%%   words  -- trees whose leaves are actual words (italicised, aligned)
%%%   cat    -- trees whose leaves are categories or feature structures
\forestset{
  words/.style={for tree={s sep=7mm, l sep=5mm,
                          if n children=0{font=\itshape, tier=terminal}{}}},
  cat/.style={for tree={s sep=7mm, l sep=5mm}},
}
\newcommand{\wtree}[1]{\begin{forest} words #1 \end{forest}}
\newcommand{\ctree}[1]{\begin{forest} cat #1 \end{forest}}
%%% the rule being applied at this step of a derivation
\newcommand{\step}[1]{\begin{center}\emp{#1}\end{center}}

%%% AVMs (Language Science Press)
\usepackage{langsci-avm}

%%% Macros the decks use, carried over from headx.tex
\newcommand{\blank}{\rule{3em}{1pt}\xspace}
\newcommand{\ft}[1]{\textsc{#1}}
\newcommand{\idn}{\mtciteform}    %%% Indonesian
\newcommand{\prd}[1]{\textbf{#1}}
\newcommand{\typ}[1]{\textit{#1}}
%%% \val: provide before renewing, in case a package defines it.
\providecommand{\val}[1]{}
\renewcommand{\val}[1]{\textit{#1}}
%%% Manning's \avmbox, still used in tree labels and in maths where a
%%% bare langsci \tag would not work.  \avm{\tag{...}} gives the same
%%% boxed tag as the ones inside feature structures.
\newcommand{\avmbox}[1]{\mbox{\avm{\tag{#1}}}}

%%% Ring a value in a dotted oval.  These used pstricks \ovalnode,
%%% which came in with rtrees; tikz (via forest) does the same job.
\newcommand{\OV}[1]{\tikz[baseline]{\node[draw,dotted,red,ellipse,
          inner sep=1pt,outer sep=0pt]{#1};}}
\newcommand{\OVB}[1]{\tikz[baseline]{\node[draw,dotted,blue,ellipse,
          inner sep=1pt,outer sep=0pt]{#1};}}

%%% From CSLI book
\newcommand{\mc}{\multicolumn}
\newcommand{\HD}{\textbf{H}\xspace}
\newcommand{\el}{< >}       %%% empty list; langsci-avm needs the space

\newcommand{\A}{\noindent\textbf{A}: }
\newcommand{\Q}{\noindent\textbf{Q}: }

%%% Question and answer on two pages, written once:
%%%   \qaslide{Title}{... question ... \answer{... answer ...}}
%%% prints the slide twice, first without the \answer parts and then
%%% with them, so the answer appears "on click".  The example counter
%%% is reset so both pages show the same example numbers.
\newif\ifshowanswer
\newcommand{\answer}[1]{\ifshowanswer#1\fi}
\newcounter{qaxnumi}
\newcommand{\qaslide}[2]{%
  \setcounter{qaxnumi}{\value{xnumi}}%
  \showanswerfalse\myslide{#1}#2%
  \setcounter{xnumi}{\value{qaxnumi}}%
  \showanswertrue\myslide{#1}#2%
  \showanswerfalse}
%\newcommand{\C}{\noindent\textbf{C}: }


\usepackage{forest}
%\avmfont{\sc}
\begin{document}
%\avmfont{\it}

\header{Lecture 3}{Feature Structures}{Headed Rules, Trees}
\maketitle



\myslide{Overview}
\begin{itemize}
\item Review: problems with CFG
\item Modeling
\item Feature structures, unification (pizza)
\item Features for linguistic description
\item Reformulate grammar rules
\item Notion of head/headedness
\item Licensing of trees
\item Reading the grammar fragment
\end{itemize}


\myslide{Our Goals}
\begin{itemize}
\item A descriptive, generative grammar
  \begin{itemize}
  \item Describing English (in this case)
  \item Generating all possible well-formed
    sentences (and no ill-formed ones)
  \item Assigning appropriate structures
  \end{itemize}
\item Design/discover an appropriate \txx{type} of
model (through incremental improvement)

\item Create a particular model (grammar
fragment) for English
\end{itemize}

\myslide{Problems with Context-Free Grammar}
\MyLogo{atomic node labels}

\begin{itemize}
\item Potentially arbitrary rules
\item Gets clunky quickly with cross-cutting properties
\item Not quite powerful enough for natural languages
\item \emp{Solution}: Replace atomic node labels (N, NP, \ldots{})  with meaningful feature structures.
\end{itemize}

\myslide{Cross-cutting Grammatical Properties}
\MyLogo{}
\vfill
\begin{large}
  \begin{tabular}[t]{l|l|l|}
    \mc{1}{c}{} & \mc{1}{c}{3rd singular subject} & \mc{1}{c}{plural subject}\\ \cline{2-3}
    direct object NP & {\it denies} & {\it deny} \\ \cline{2-3}
    no direct object NP & {\it disappears} & {\it disappear} \\ \cline{2-3}
  \end{tabular}
\end{large}

\myslide{Language Models}

  \begin{itemize}
  \item Two kinds of language models
    \begin{itemize}
    \item Speakers’ internalized knowledge (their grammar)
    \item Set of sentences in the language
    \item Probability distribution over strings \\ (what ``language
      model'' usually means today: it scores sentences, but doesn't
      license or rule them out)
    \end{itemize}
  \item Things involved in modelling language
    \begin{itemize}
    \item Real world entities (utterance types)
    \item Models (fully specified trees)
    \item Descriptions of the models (rules, principles, lexical entries)
    \end{itemize}
  \end{itemize}

\section{A delicious introduction to Typed Feature Structures}

  \myslide{Feature Structure Descriptions}
\vfill
\begin{Large}
  \avm{[ feature$_1$ &\ \ \  value$_1$\\
    feature$_2$ &\ \ \ value$_2$\\
    \hfil.\ .\ .\hfil\\
    feature$_n$ &\ \ \  value$_n$]}
\end{Large}

\bigskip
\bigskip
\bigskip

Like a form: each field (feature) has an entry (value).

Or, if you program: like a dictionary in Python (feature = key).

\myslide{A Pizza Type Hierarchy}
%\vspace*{-.5ex}
\begin{forest}
cat
[{\val{pizza-thing}} [{\avm{[\type*{pizza} 
      crust & \{ \normalfont thick, thin, stuffed \}
     \\ toppings & topping-set ]}}]
[{\avm{[\type*{topping-set} olives & \{ +, -\} 
      \\ onions & \{ +, -\} \\ mushrooms& \{ +, -\}  ]}} [{\val{vegetarian}}]
[{\avm{[\type*{non-vegetarian} sausage & \{ +, -\} 
         \\ pepperoni & \{ +, -\} \\ ham & \{ +, -\} ]}}]]]
\end{forest}
\vspace*{-2.5ex}

\myslide{Types}

\begin{tabular}{lll}
  Type & Features/Values & IST \\ \hline
  \smaller[2]\it pizza-thing & & NONE\\
  \smaller[2]\it  pizza & {\small  \avm{%
    [ crust & \{ \normalfont thick, thin, stuffed\}
    \\ toppings & \it topping-set ]}}   & \smaller[2]\it  pizza-thing \\
  \smaller[2]\it topping-set &{\small  \avm{%
      [  olives & \{ +, -\} 
    \\ onions & \{ +, -\} \\ mushrooms& \{ +, -\}  ]}}   & \smaller[2]\it pizza-thing \\
  \smaller[2]\it vegetarian &  & \smaller[2]\it topping-set \\
 \smaller[2]\it non-vegetarian &{\small  \avm{%%%
     [ sausage & \{ +, -\}  \\ pepperoni & \{ +, -\} 
     \\ ham & \{ +, -\} ]}}   & \smaller[2]\it topping-set \\

\end{tabular}

\txx{IST}: Immediate Super Type

\myslide{Type Hierarchies}

A type hierarchy \ldots
\begin{itemize}
\item[\ldots]  states what kinds of objects we claim exist\\
  --- the types

\item[\ldots]  organizes the objects hierarchically into classes
  with shared properties \\
  --- the type hierarchy

\item[\ldots]   states what general properties each kind of object
has \\ --- the feature and feature value declarations
\end{itemize}

\qaslide{Pizza Descriptions and Pizza Models}{
\begin{center}
  \avm{%
    [\type*{pizza}
    crust & thick \\
    toppings &
    [\type*{vegetarian} olives & + \\ onions & + ] ]}
\end{center}
How many pizza models (by definition, fully resolved) satisfy this description?
\answer{\par
\emp{Two}: MUSHROOMS must have a value, and it can be $+$ or $-$.

\begin{center}
  \avm{%
  [\type*{pizza}
  crust & thick \\
  toppings &
  [\type*{vegetarian} olives & + \\ onions & + \\ mushrooms& $-$ ] ]}
\avm{%
  [\type*{pizza}
  crust & thick \\
  toppings &
  [\type*{vegetarian} olives & + \\ onions & + \\ mushrooms & + ] ]}
\end{center}}
}

\qaslide{Pizzas in the World}{
\begin{center}
  \avm{%
    [\type*{pizza}
    crust & thick \\
    toppings &
    [\type*{vegetarian} olives & + \\ onions & + ] ]}
\end{center}
\begin{itemize}
\item How many pizzas-in-the-world do the pizza
models correspond to?
\answer{\\ \emp{A large, constantly-changing number.}
\item The ‘type’/‘token’ distinction applies to sentences as well}
\end{itemize}
}

\qaslide{Combining Constraints}{
\begin{center}
  \avm{%
    [\type*{pizza}
    crust & thick \\
    toppings &
    [ olives & + \\ ham & $-$ ] ]
    \& 
   [\type*{pizza}
    toppings &
    [ olives & + \\ onions & + ] ]}
  \answer{\\[2ex] = \
  \avm{%
    [\type*{pizza}
    crust & thick \\
    toppings &
    [ olives & + \\ ham & $-$ \\ onions & + ] ]}}
\end{center}

Combining constraints is called \txx{unification}, also written
$\sqcap$.  Information is added, never lost.
}



\qaslide{Combining Constraints}{
\begin{center}
  \avm{%
    [\type*{pizza}
    crust & thick \\
    toppings &
    [ olives & + \\ ham & $-$ ] ]
    \& 
   [\type*{pizza}
   crust & thin \\
    toppings &
    [ olives & + \\ onions & + ] ]} 
\answer{\begin{huge}
\\[2ex]   = $\varphi$
\end{huge}}
\end{center}

\answer{\emp{A crust can't be both thick and thin}, so unification
fails.  Failure (NULL) is also written as:  $\perp,  \emptyset, \phi$.}
}

%\myslide{Combining Constraints}

\qaslide{Combining Constraints}{
\begin{center}
  \avm{%
    [\type*{pizza}
    crust & thin \\
    toppings &
    [ olives & + \\ ham & $-$ ] ]
    \& 
   [\type*{pizza}
   crust & thin \\
    toppings &
    [\type*{vegetarian} ] ]} 
\answer{\begin{huge}
\\[2ex]   = $\varphi$
\end{huge}}
\end{center}

\answer{\emp{\textit{vegetarian} has no feature \ft{ham}!}  (Even
though the value is $-$.)}
}

\myslide{A New Theory of Pizzas}
\begin{center}
  \avm{%
    [\type*{pizza}
    crust & \{ \normalfont thick, thin, stuffed \} \\
    one-half & \val{topping-set} \\
    other-half & \val{topping-set} ]}
\end{center}


\qaslide{Combining Constraints}{
\begin{center}
  \avm{%
    [\type*{pizza}
    one-half &
    [ olives & + \\ ham & $-$ ] ]
    \& 
   [\type*{pizza}
    other-half &
    [  olives & $-$ \\ 
         ham & + ] ]}
  \answer{\\ = 
  \avm{%
 [\type*{pizza}
   one-half &
   [ olives & + \\ ham & $-$ ] \\
   other-half &
   [  olives & $-$ \\ ham & + ] ]}}
\end{center}
}


\myslide{Identity Constraints (tags)}
\begin{center}
   \avm{%
 [\type*{pizza}
   one-half &
   [ olives & \1 \\ ham & \2 ] \\
   other-half &
   [  olives & \1 \\ ham & \2 ] ]} 
\end{center}

A boxed number (a \txx{tag}) means \emp{the very same value}:
whatever \ft{olives} is on one half, it is on the other half too.

\qaslide{Combining Constraints}{
\begin{center}
  \avm{%
    [\type*{pizza}
    one-half &
    \1 [ olives & + \\ ham & $-$ ]
    \\   other-half & \1 ]
    \& 
   [\type*{pizza}
    other-half &
    [  olives & + \\ mushrooms & + ] ]}
  \answer{\\ = 
  \avm{%
 [\type*{pizza}
   one-half &
    \1 [ olives & + \\ ham & $-$ \\ mushrooms & + ] \\
   other-half & \1 ]}}
\end{center}

\answer{\emp{Mushrooms on the other half means mushrooms on both halves.}}
}


\myslide{Note}
\begin{center}
 \avm{%
 [\type*{pizza}
   one-half &
    \1 [ olives & + \\ ham & $-$ \\ mushrooms & + ] \\
   other-half & \1 ]} 
\\ =
  \avm{%
 [\type*{pizza}
   one-half &
    \1 \\
   other-half & \1 [ olives & + \\ ham & $-$ \\ mushrooms & + ] ]} 
\end{center}

It doesn't matter where we write the value: both paths lead to the
same thing.

\qaslide{Combining Constraints}{
\begin{center}
  \avm{%
    [\type*{pizza}
    crust & thick \\
    one-half &
    \1 [ olives & + \\ onions & $-$ ] 
    \\ other-half & \1  \textit{vegetarian} ]
    \& 
   [\type*{pizza}
   crust & thin \\
    one-half &
    [ olives & + \\ pepperoni & + ] ]} 
\begin{huge}
\\[2ex]   = $\varphi$
\end{huge}
\end{center}

Why does this fail?  (Two reasons.)
\answer{\par
\emp{(1) the crust can't be both thick and thin; (2) one half is the
same as the other half, which is \textit{vegetarian}, and
\textit{vegetarian} has no \ft{pepperoni}.}}
}



\qaslide{Why combine constraints?}{
\begin{itemize}
\item The pizza example illustrates how
unification can be used to combine
information from different sources.

\item In our grammar, information will come
from lexical entries, grammar rules, and
general principles.
\item What makes us put the information together?
\answer{\\ \emp{The definition of a well-formed tree: every node must
    satisfy all of them at once.}}
\end{itemize}
}

\section{Describing Language with TFS}
  

\myslide{Linguistic Application of Feature Structures:}
Making the Mnemonic Meaningful
\begin{itemize}
\item What do these CFG categories have in common?
  \begin{itemize}
  \item  NP \& VP:  are both phrases
  \item N \& V:  are both words
  \item NP \& N: are both ‘nouny’
  \item VP \& V:  are both ‘verby’
  \end{itemize}
\end{itemize}

\myslide{The Beginnings of Our Type Hierarchy}


\begin{center}
  \begin{forest}
cat, for tree={font=\itshape}
[{feat-struc} [{expression} [{word}]
[{phrase}]]
[{pos} [{noun}]
[{verb}]
[{det}]
[{prep}]
[{adj}]
[{conj}]]]
\end{forest}
\end{center}
\vspace*{5ex}

  \begin{tabular}{lclc}
NP &  \avm{%
[\type*{phrase}
   head & \type{noun} ]} 
& VP & 
  \avm{%
[\type*{phrase}
   head & \type{verb} ]} \\
N &
\avm{%
< \textnormal{fly} , \ [\type*{word}
               head & \type{noun} ]>}
& V &
\avm{%
< \textnormal{fly} , \ [\type*{word}
               head & \type{verb} ]>}
  \end{tabular}




\myslide{Type Hierarchy for Parts of Speech II}

\begin{center}
  \begin{forest}
cat, for tree={font=\itshape}
[{feat-struc} [{expression  \avm{[ \textsc{head} ]}} [{word}]
[{phrase}]]
[{pos} [{agr-pos \avm{[ \textsc{agr} ]}} [{noun}]
[{verb \avm{[ \textsc{aux} ]}}]
[{det}]]
[{prep}]
[{adj}]
[{conj}]]]
\end{forest}
\end{center}


\myslide{Agreement}

We need more information to make words agree.

\avm{%
< \textnormal{fly}, \ [\type*{word}
               head & [\type*{noun}
                       agr & [\type*{agr-cat} per & 3rd \\ num & sg ] ] ]>}

\avm{%
< \textnormal{flies}, \ [\type*{word}
               head & [\type*{verb}
                       agr & [\type*{agr-cat} per & 3rd \\ num & sg ] ] ]>}

\myslide{Agreement (S  $\rightarrow$  NP VP)}

\avm{%
[\type*{phrase}\\
  head & {\1}\type{verb}\\
  val & [ comps & itr\\
           spr & $+$]] \ $\rightarrow$\ \ [\hfil NP\\
                              [ head \ [agr & {\2}]]] \ \ [\type*{phrase}\\
                                        head & {\1}[agr & {\2}]\\
                                        val & [ comps & itr\\
                                                 spr & $-$]]}

The values on \ft{agr} for the subject NP and verb phrase must be identical.


\myslide{A Simple Feature for Valence}

\begin{tabular}[t]{cc}
\avm{%
  iv =  [\type*{word}\\
  head & \type{verb}\\
  val & [\type*{val-cat}\\
  comps & itr]]} &
\avm{%
  tv  =  [\type*{word}\\
  head & \type{verb}\\
  val & [\type*{val-cat}\\
  comps & str ]]}\\
\avm{%
  dtv =  [\type*{word}\\
  head & \type{verb}\\
  val & [\type*{val-cat}\\
  comps & dtr ]]}
\end{tabular}

\ft{comps} controls how many \txx{complements} are possible.


\myslide{Head-Complement Rules (VP $\rightarrow$ V $|$ V NP $|$ V NP NP)}
\begin{small}
1 \avm{% 
[\type*{phrase}\\
  head & {\1} \\
  val & [ comps & itr\\
           spr & $-$]] \ $\rightarrow$\ \ [\type*{word}\\
                                               head & {\1}\\
                                               val & [ comps & itr\\
                                                        spr & $-$]]}

2 \avm{% 
[\type*{phrase}\\
  head & {\1}\\
  val & [ comps & itr\\
           spr & $-$ ]] \ $\rightarrow$\ \ [\type*{word}\\
                                               head & {\1}\\
                                               val & [ comps & str\\
                                                        spr & $-$]]\ \ NP}

3 \avm{%
[\type*{phrase}\\
  head & {\1}\\
  val & [ comps & itr\\
           spr & $-$ ]] \ $\rightarrow$\ \ [\type*{word}\\
                                               head & {\1}\\
                                               val & [ comps & dtr\\
                                                        spr & $-$]]\ \ NP\ \ NP}
\end{small}

\myslide{Underspecification}

\avm{%
 V \ = \ [\type*{word}\\
              head & \type{verb}]}

\avm{% 
            VP \ = \ [\type*{phrase}\\
              head & \type{verb}]}

\avm{%
 [ head & \type{verb}]} 

This captures \textit{verbiness}!

\myslide{Another Valence Feature}
\begin{tabular}[t]{cc}
\avm{\ \ \ \ NP  =  \ [\type*{phrase}\\
                              head & \type{noun} \\
			      val & [\type*{val-cat}\\
				       comps & itr\\
				       spr & $+$]]}&
\avm{NOM =  \ [\type*{phrase}\\
                head & \type{noun} \\
		val & [\type*{val-cat}\\
			 comps & itr\\
                         spr & $-$]]}
\end{tabular}


\ft{spr} controls the \txx{specifier} (determiner and/or subject)

\begin{center}
\begin{forest}
  [ NP [ NOM [ N ]]]
\end{forest}
\end{center}


\myslide{Specifier and Verbs}

\begin{tabular}[t]{cc}
\avm{\hskip 11pt\ \ \ S  =  \ [\type*{phrase}\\
                                       head & \type{verb} \\
				       val & [\type*{val-cat}\\
					       comps & itr\\
					       spr & $+$]]}&
\avm{\ \ VP =  \ [\type*{phrase}\\
                head & \type{verb} \\  
	        val & [\type*{val-cat}\\
			 comps & itr\\
                         spr & $-$]]}
\end{tabular}
\begin{center}
\begin{forest}
  [ S [ VP [ V ]]]
\end{forest}
\end{center}

S is the verbal counterpart of NP, and VP the verbal counterpart of NOM.

\myslide{S and NP}

\avm{[ val & [\type*{val-cat}\\
	    comps & itr\\
            spr & $+$ ]]}

\begin{itemize}
\item both are fully saturated: specified and no more complements
  \begin{exe}
    \ex \eng{We created a monster}
    \ex \eng{our creation of a monster}
  \end{exe}
\end{itemize}

\myslide{Type Hierarchy So Far}

\begin{center}\small
  \scalebox{0.9}{\begin{forest}
cat, for tree={font=\itshape}
[{feat-struc} [{expression  \avm{[ \textsc{head, val} ]}} [{word}]
[{phrase}]]
[{pos} [{agr-pos \avm{[ \textsc{agr} ]}} [{noun}]
[{verb \avm{[ \textsc{aux} ]}}]
[{det}]]
[{prep}]
[{adj}]
[{conj}]]
[{val-cat \avm{[ \textsc{spr, comps} ]}}]
[{agr-cat \avm{[ \textsc{per, num} ]}}]]
\end{forest}}
\end{center}


\qaslide{Heads}{
\begin{itemize}
\item Intuitive idea: A phrase typically contains a word that
  determines its most essential properties, including
  \begin{itemize}
  \item where it occurs in larger phrases
  \item what its internal structure is
  \end{itemize}
\item This is called the head
\item The term \txx{head} is used both for the head word in a
  phrase and for all the intermediate phrases containing
  that word
\item NB: Not all phrases have heads
\\ can you think of a phrase that doesn't?
\answer{\\ \emp{Coordination: \lex{Kim and Sandy}}}
\end{itemize}
}

\myslide{Formalizing the Notion of Head}
\begin{itemize}
\item Expressions have a feature \ft{head}
\item \ft{head}’s values are of type \val{pos} (part-of-speech)
\item For \ft{head} values of type \val{agr-pos}, \ft{head}’s
  value also includes the feature \ft{agr}
\item Well-formed trees are subject to the \txx{Head Feature Principle}
\end{itemize}

\myslide{The Head Feature Principle}
\begin{itemize}
\item Intuitive idea: Key properties of phrases are
shared with their heads
\item The \txx{HFP}:
  \begin{quote} \large    In any headed phrase, the \ft{head}
    value of the mother and the head daughter
    must be identical.    
  \end{quote}
\item Sometimes described in terms of properties
“percolating up” or “filtering down”, but this
is just metaphorical talk
\item  the \index{head daughter (H)} head daughter in a headed-rule will be labeled with  `{\bf H}'.

  \begin{quote}
    [ \typ{type} ]  \ \ \into  \ \ \ldots \ \ \HD [\ \  \ ]  \ \ \ldots
  \end{quote}


\end{itemize}

\myslide{A Tree is Well-Formed if \ldots }
\begin{itemize}
\item It and each subtree are licensed by a grammar rule
or lexical entry
\item All general principles (like the HFP) are satisfied.
\item NB: Trees are part of our model of the language,
so all their features have values (even though we
will often be economical and leave out the values
irrelevant to our current point).
\end{itemize}

\qaslide{Question: Unheaded Phrases}{
Do phrases that are not headed have
\ft{head} features?
\answer{\par
\bigskip
\emp{Yes}, all expressions have \ft{head}, but in a non-headed
phrase it does not have to be identical to any daughter's.}
}

\myslide{Reformulating the Grammar Rules I}

Which simple phrase structure rules (Ch 2) do these correspond to?

\begin{itemize}
\item  Head-Complement Rule 1:

\avm{[\type*{phrase}\\
  val & [ comps & itr\\
           spr & $-$]] \ $\rightarrow$\ \HD \ [\type*{word}\\
                                               val & [ comps & itr\\
                                                        spr & $-$]]}
\item  Head-Complement Rule 2:

\avm{[\type*{phrase}\\
  val & [ comps & itr\\
           spr & $-$ ]] \ $\rightarrow$\ \HD \ [\type*{word}\\
                                               val & [ comps & str\\
                                                        spr & $-$]]\ \ NP}
\end{itemize}

\myslide{Reformulating the Grammar Rules I (cont.)}
\begin{itemize}
\item  Head-Complement Rule 3:

\avm{[\type*{phrase}\\
  val & [ comps & itr\\
           spr & $-$ ]] \ $\rightarrow$\ \HD \ [\type*{word}\\
                                               val & [ comps & dtr\\
                                                        spr & $-$]]\ \ NP\ \ NP}
\end{itemize}

\myslide{Reformulating the Grammar Rules II}
\MyLogo{Head-Specifier Rule 2 updated in the problems}

\begin{itemize}
\item Head-Specifier Rule 1:

\avm{[\type*{phrase}\\
val & [ comps & itr\\
         spr & $+$]] \ $\rightarrow$\ \ [\hfil NP\\
[head \ [agr \ {\1}]]]\ \ \ {\HD}[\type*{phrase}\\
                                        head & [\type*{verb}\\
agr & {\1}]\\
                                        val & [ spr & $-$]]}
\item Head-Specifier Rule 2:

\avm{[\type*{phrase}\\
val & [ comps & itr\\
         spr & $+$]] \ $\rightarrow$\ \ D \ \ {\HD}[\type*{phrase}\\
                                        head & \type{noun}\\
                                        val & [ spr & $-$]]}
\end{itemize}

\myslide{Reformulating the Grammar Rules III}
\begin{itemize}
\item Non-Branching NP Rule

\avm{[\type*{phrase}\\
%head & \type{noun}\\
val & [ comps & itr\\
         spr & $+$]] \ $\rightarrow$\ \ {\HD}[\type*{word}\\
                                     head & \type{noun}\\
                                     val & [ spr & $+$]]}

\item Head-Modifier Rule

\avm{[\type*{phrase}\\
		    val & [ comps & itr\\
		             spr & $-$ ]] \ 
$\rightarrow$\ \ {\HD}[\type*{phrase}\\
			     val & [ spr & $-$]] \
                 PP}
\end{itemize}

\myslide{Reformulating the Grammar Rules III (cont.)}
\begin{itemize}
\item Coordination Rule

\avm{ [ head \1 ]\ \ \ $\rightarrow$\ \ \ %
[ head \1 ]+\ [\type*{word}\\
 head & \type{conj} ]\  [ head \1 ]}
\\[4ex]
Only coordinate two similar things!
\end{itemize}

\qaslide{Which Chapter 2 Rules Do the New Rules Replace?}{
\begin{center}
  \renewcommand{\arraystretch}{1.3}
  \begin{tabular}{ll}
    New rule & Chapter 2 rules \\ \hline
    Head-Complement Rule 1 & \answer{VP $\rightarrow$ V \quad NOM $\rightarrow$ N} \\
    Head-Complement Rule 2 & \answer{VP $\rightarrow$ V NP \quad \emp{and} PP $\rightarrow$ P NP} \\
    Head-Complement Rule 3 & \answer{VP $\rightarrow$ V NP NP} \\
    Head-Specifier Rule 1 & \answer{S $\rightarrow$ NP VP} \\
    Head-Specifier Rule 2 & \answer{NP $\rightarrow$ D NOM \quad (D now \emp{required})} \\
    Non-Branching NP Rule & \answer{NP $\rightarrow$ NOM, with \emp{no} D (pronouns, names)} \\
    Head-Modifier Rule & \answer{VP $\rightarrow$ VP PP \quad NOM $\rightarrow$ NOM PP} \\
    Coordination Rule & \answer{X $\rightarrow$ X$^+$ CONJ X} \\
  \end{tabular}
\end{center}
\answer{It is not one-to-one: one new rule can do the work of several old ones.}
}

\qaslide{Advantages of the New Formulation}{
\begin{itemize}
\item Subject-verb agreement is stipulated only
once (where?)
\answer{\\ \emp{in Head-Specifier Rule 1}}
\item Common properties of verbs with different
valences are expressed by common features
(for example?)
\answer{\\ \emp{\ft{agr} (\ft{per}, \ft{num}): \lex{denies} and
    \lex{disappears} are both [\ft{per} 3rd, \ft{num} sg]}}
\item Parallelisms across phrase types are captured
(for example?)
\answer{\\ \emp{Head-Complement Rule 2 builds both VPs and PPs}}
\end{itemize}
}

\qaslide{Disadvantages of the New Formulation}{
\begin{itemize}
\item We still have three head complement rules
\item We still have two head specifier rules
\item We only deal with three verb valences
\\ (Which ones? What are some others?)
\answer{\\ \emp{intransitive, transitive, ditransitive; but also
    \lex{put} (NP PP), \lex{think} (S), \lex{tell} (NP S), \ldots}}
\item The non-branching rule doesn’t really do any
empirical work
\item Anything else?
\end{itemize}
}



%%% They swim, built up step by step.  \theytree's five arguments are
%%% the frame colours of the S, NP, they, VP and swim nodes (white =
%%% invisible), so every page has the same layout.  #6 is the text on
%%% the right.
\newcommand{\nd}[2]{{\setlength{\fboxsep}{1pt}\fcolorbox{#1}{white}{#2}}}
\newcommand{\theytree}[6]{%
\noindent\begin{minipage}[t]{0.6\textwidth}\vspace{0pt}
\scalebox{0.6}{\begin{forest}
cat, for tree={l sep=4mm}
[{\nd{#1}{\avm{[\type*{phrase}\\
  head &[\type*{verb}\\
  agr &[  per & 3rd\\
  num & pl]]\\
  val & [  comps & itr\\
  spr & $+$ ]]}}}
 [{\nd{#2}{\avm{[\type*{phrase}\\
  head & [\type*{noun}\\
  agr & [  per & 3rd\\
  num & pl]]\\
  val & [  comps & itr\\
  spr & $+$ ]]}}}
  [{\nd{#3}{\avm{[\type*{word}\\
  head &[\type*{noun}\\
  agr & [  per & 3rd\\
  num & pl]]\\
  val & [  comps & itr\\
  spr & $+$ ]]}}} [{they}]]]
 [{\nd{#4}{\avm{[\type*{phrase}\\
  head &[\type*{verb}\\
  agr &[  per & 3rd\\
  num & pl]]\\
  val & [  comps & itr\\
  spr & $-$ ]]}}}
  [{\nd{#5}{\avm{[\type*{word}\\
  head &[\type*{verb}\\
  agr &[  per & 3rd\\
  num & pl]]\\
  val & [  comps & itr\\
  spr & $-$ ]]}}} [{swim}]]]]
\end{forest}}
\end{minipage}\hfill
\begin{minipage}[t]{0.38\textwidth}\vspace{0pt}\small
#6
\end{minipage}}

\myslide{Which Rule Licenses Each Node?}
\theytree{white}{white}{white}{white}{white}{%
  \textbf{Which rule or lexical entry licenses each node?}}

\myslide{Which Rule Licenses Each Node?}
\theytree{white}{white}{red}{white}{red}{%
  \textbf{Which rule or lexical entry licenses each node?}
  \begin{itemize}
  \item \lex{they}, \lex{swim}: their \emp{lexical entries}
  \end{itemize}}

\myslide{Which Rule Licenses Each Node?}
\theytree{white}{red}{white}{white}{white}{%
  \textbf{Which rule or lexical entry licenses each node?}
  \begin{itemize}
  \item \lex{they}, \lex{swim}: their lexical entries
  \item NP: the \emp{Non-Branching NP Rule} (\lex{they} is
    [\ft{spr} $+$]); HFP
  \end{itemize}}

\myslide{Which Rule Licenses Each Node?}
\theytree{white}{white}{white}{red}{white}{%
  \textbf{Which rule or lexical entry licenses each node?}
  \begin{itemize}
  \item \lex{they}, \lex{swim}: their lexical entries
  \item NP: the Non-Branching NP Rule (\lex{they} is
    [\ft{spr} $+$]); HFP
  \item VP: \emp{Head-Complement Rule 1} (\lex{swim} is
    [\ft{comps} \val{itr}]); HFP
  \end{itemize}}

\myslide{Which Rule Licenses Each Node?}
\theytree{red}{white}{white}{white}{white}{%
  \textbf{Which rule or lexical entry licenses each node?}
  \begin{itemize}
  \item \lex{they}, \lex{swim}: their lexical entries
  \item NP: the Non-Branching NP Rule (\lex{they} is
    [\ft{spr} $+$]); HFP
  \item VP: Head-Complement Rule 1 (\lex{swim} is
    [\ft{comps} \val{itr}]); HFP
  \item S: \emp{Head-Specifier Rule 1}, which also makes the NP's
    and VP's \ft{agr} the same; HFP
  \end{itemize}}

\myslide{In abbreviated form}

\begin{forest}
cat
[{S   \avm{[ \ft{head} & \1 [ \ft{agr} & \4
    [ \ft{per} & 3 \\ \ft{num} & pl ] ] ]}} [{NP  \avm{[ \ft{head} \2 [ \ft{agr} \4 ] ]}} [{N  \avm{[ \ft{head} \2 ]}} [{\textnormal{they}}]]]
[{VP  \avm{[ \ft{head} \1 ]}} [{V  \avm{[ \ft{head} \1 ]}} [{\textnormal{swim}}]]]]
\end{forest}
\vspace*{-2ex}
\begin{flushleft}
  % HFP $\Rightarrow$ head the same for mother and daughter; \\
  S $\Rightarrow$ \val{phrase}, \ft{head} \val{verb}, \ft{comps}
  \val{itr}, \ft{spr} \val{+};  \\
  VP $\Rightarrow$ \val{phrase}, \ft{head}
  \val{verb}, \ft{comps} \val{itr}, \ft{spr} \val{$-$};   \ldots
\end{flushleft}
\myslide{A Question:}
Since the lexical entry for swim below has only [NUM pl] as
the value of \ft{agr}, how did the tree on the previous slide get
[PER 3rd] in the \ft{agr} of swim?

\avm{%
< \textnormal{swim}, \ 
[\type*{word} 
  head & [\type*{verb}  agr &   [  num & pl ] ] \\
            val &  [\type*{val-cat}  comps & itr \\ spr & $-$ ] ] 
 >}


\myslide{Answer: A Chain of Identities}
\begin{center}
  \renewcommand{\arraystretch}{1.3}
  \begin{tabular}{lll}
    This \ft{agr} value \ldots & is identical to \ldots & because of \ldots \\ \hline
    \lex{they}'s lexical entry & \lex{they}'s word node & the lexical entry \\
    \lex{they}'s word node & the NP & the HFP \\
    the NP & the VP & Head-Specifier Rule 1 \\
    the VP & \lex{swim}'s word node & the HFP \\
  \end{tabular}
\end{center}
So \lex{swim}'s \ft{agr} is the very same value as \lex{they}'s, which
is [\ft{per} 3rd, \ft{num} pl].  \emp{Constraints from different
places combine}, just like the pizza.

\section{Reading the Grammar Fragment}

\myslide{Components of the Grammar}
The grammar summary at the end of each chapter (and in the Appendix)
has four parts:
\begin{enumerate}
\item The \txx{type hierarchy}, with the features and values
  appropriate for each type
\item The \txx{phrase structure rules}
\item The \txx{well-formedness conditions} / general principles
  (e.g.\ the HFP)
\item The \txx{lexical entries}
\end{enumerate}
A tree is well-formed if it satisfies all of them.

\myslide{What about the Abbreviations?}
\begin{itemize}
\item NP, VP, S, N, V \ldots\ are just for convenience
\item Each abbreviation is defined as an underspecified feature
  structure, e.g.
  \begin{center}
    \avm{NP = [\type*{phrase}
      head & \type{noun} \\
      val & [ comps & itr \\ spr & $+$ ] ]}
  \end{center}
\item We use them in trees as shorthand, but they do no work in the
  grammar
\end{itemize}

\myslide{Building Trees}
\begin{itemize}
\item You should always be able to check your trees against the grammar
\item The grammar should give you enough information to fill in every
  feature on every node
\item The book only lists a few lexical entries: find the closest one
  to your word and use it as a guide
\item Sometimes there is a ``bug'' in the grammar in the book: if you
  can't work out a value from the grammar summary and the principles,
  \emp{ask}!
\end{itemize}

\myslide{Summary}
\begin{itemize}
\item Review: problems with CFG; modeling
\item Feature structures, unification (pizza)
\item Features for linguistic description
\item Reformulate grammar rules
\item Notion of head/headedness
\item Licensing of trees
\item Reading the grammar fragment
\item Next time: Valence and agreement: complex feature values
\end{itemize}




\myslide{3.1 Applying the grammar}
\begin{itemize}
\item[A.] Formulate \index{agreement|(} a \index{lexical entry} lexical entry for the word {\it defendants}.
\item[B.] Draw a tree for the sentence {\it The defendants walk}.  Show
the values for all of the features on every node and use tags to indicate
the effects of any identities that the %rules require.
grammar requires.
\item[C.] Explain how your lexical entry for {\it defendants} interacts
with the Chapter 3 grammar to rule out {\it *The defendants walks}.
Your explanation should make reference to grammar rules, lexical
entries and the \index{Head Feature Principle (HFP)}HFP. \index{agreement|)}
\end{itemize}

\myslide{3.3 Determiner-Noun Agreement}

The Chapter 3 grammar declares AGR to be a feature
appropriate for the types {\it noun}, {\it verb}, and {\it det}, but
so far we haven't discussed agreement involving determiners.  Unlike
the determiner {\it the}, most other English determiners do show
agreement with the nouns they combine with:

\begin{itemize}\addtolength{\itemsep}{-1ex}
\item[(i)] a bird/*a birds
\item[(ii)] this bird/*this birds
\item[(iii)] that bird/*that birds
\item[(iv)] these birds/*these bird
\item[(v)] those birds/*those bird
\item[(vi)] many birds/*many bird
\end{itemize}
\newpage

\begin{itemize}
\item[A.] Formulate lexical entries for {\it this} and {\it these}.
\item[B.] Modify Head-Specifier Rule 2 so that it enforces agreement
between the noun and the determiner just like Head-Specifier Rule 1
enforces agreement between the NP and the VP.
\item[C.] Draw a tree for the NP {\it these birds}. Show the
value for all features of every node and use tags to indicate the
effects of any identities that the %rules 
grammar 
(including your modified
HSR2 and the Head Feature Principle) requires.
\end{itemize}


\myslide{Types for English Agreement}

\begin{forest}
  /tikz/every node/.append style={font=\it},
  [agr-cat [sg,name=sg [3sg]] [, phantom] [, phantom]  [, phantom] [non-3sg
  [1sg, name=1sg] [non-1sg [2sg, name=2sg] [pl]]]]
  \draw (sg.south) -- (1sg.north);
  \draw (sg.south) -- (2sg.north);
\end{forest}

\myslide{Acknowledgments and References}
\begin{itemize}
\item Course design and slides borrow heavily from Emily Bender's course:
\textit{Linguistics 566: Introduction to Syntax for Computational Linguistics}
\\ \url{http://courses.washington.edu/ling566}
\item This lecture complements \citet[Chapter 3]{sag:wasow:bender:2003},
  and the problems are taken from it
\end{itemize}




\bibliographystyle{aclnat}
\bibliography{abb,mtg,nlp,ling}
\end{document}


%%% Local Variables: 
%%% coding: utf-8
%%% mode: latex
%%% TeX-PDF-mode: t
%%% TeX-engine: xetex
%%% End: 
